\chapter{Syzygies}
\label{ch:v24}

Syzygies are relations among generators, then relations among relations, and so on. Iterating syzygies turns presentations into resolutions and exposes hidden structure in modules.

\section*{Learning goals}
After this chapter, the reader should be able to:
\begin{itemize}
\item compute first syzygies as relations among a chosen generating set;
\item construct higher syzygies recursively from kernels in a free resolution;
\item translate presentation matrices into syzygy modules;
\item analyze how changing generators changes a presentation but not the underlying module;
\item use syzygies to continue a free resolution;
\item distinguish minimal relations from redundant generators or redundant syzygies;
\end{itemize}
\section*{Conceptual roadmap}
\[
\boxed{\text{presentation or universal property}\;\longrightarrow\;\text{exactness/localization}\;\longrightarrow\;\text{structural consequence}.}
\]

\section{First syzygy}
Given a surjection $F_0\to M$, the first syzygy module is its kernel.

\section{Higher syzygies}
The $n$th syzygy is the kernel at stage $n-1$ of a chosen free or projective resolution.

% BEGIN VOL05-EXPANSION V24-example-01
\begin{example}[A first syzygy]\label{ex:v24-expand-01}
For the ideal $(x,y)\subset k[x,y]$, the map $R^2\to(x,y)$ sends $(a,b)$ to $ax+by$. The vector $(-y,x)$ lies in the kernel, giving the basic syzygy $-yx+xy=0$.
\end{example}
% END VOL05-EXPANSION V24-example-01
\section{Dependence on presentation}
Syzygies depend on choices, but stable projective information is controlled up to adding free summands.

\section{Presentations}
A finite presentation gives a finitely generated first syzygy.

\section{Noetherian setting}
Over a Noetherian ring, syzygies of finite modules in finite free presentations remain finitely generated.

% BEGIN VOL05-EXPANSION V24-example-02
\begin{example}[Syzygy from a presentation matrix]\label{ex:v24-expand-02}
If $R^m\xrightarrow{A}R^n\to M\to0$ presents $M$, then the first syzygy module is $\ker A$. A second syzygy is obtained by presenting that kernel and taking the next kernel.
\end{example}
% END VOL05-EXPANSION V24-example-02
\section{Dimension shifting}
Higher Tor and Ext groups can be expressed through lower groups of syzygies.

\section{Regular sequences}
Koszul-type relations provide canonical syzygies for complete-intersection situations.

% BEGIN VOL05-EXPANSION V24-example-03
\begin{example}[Koszul relation]\label{ex:v24-expand-03}
For a regular pair $x,y$, the relation $(-y,x)$ generates the first relation among the two generators of $(x,y)$. This is the opening step of the Koszul resolution.
\end{example}
% END VOL05-EXPANSION V24-example-03
\section{Hilbert syzygy preview}
Over polynomial rings over fields, finite modules have bounded projective resolution length.

\section{Computational matrices}
Syzygies are kernels of matrices and can be studied algorithmically.

\section{Core structural results}

\begin{theorem}[Syzygies over Noetherian rings are finite]\label{thm:v24-01}
If $A$ is Noetherian and $M$ is finitely generated, then the kernel of a map $A^n\to M$ is finitely generated.
\begin{proof}
The finite free module $A^n$ is Noetherian, so every submodule, including the kernel, is finitely generated.
\end{proof}
\end{theorem}

\begin{theorem}[Schanuel lemma]\label{thm:v24-02}
If $0\to K\to P\to M\to0$ and $0\to K'\to P'\to M\to0$ have projective $P,P'$, then $K\oplus P'\cong K'\oplus P$.
\begin{proof}
Form the pullback of $P\to M$ and $P'\to M$. Its two projections fit into short exact sequences that split because $P$ and $P'$ are projective, yielding the two decompositions.
\end{proof}
\end{theorem}

\begin{theorem}[Dimension shifting for Tor]\label{thm:v24-03}
If $0\to\Omega M\to P\to M\to0$ with $P$ projective, then $\operatorname{Tor}_{n}^A(M,N)\cong\operatorname{Tor}_{n-1}^A(\Omega M,N)$ for $n\ge2$.
\begin{proof}
Apply the long exact Tor sequence. Higher Tor of the projective module $P$ vanishes, so adjacent terms give the stated isomorphism.
\end{proof}
\end{theorem}

\section{Worked examples}

\begin{example}[First syzygy diagnostic]\label{ex:v24-worked-01}
Give a rigorous worked analysis of First syzygy. State the ring, module, finiteness, localization, or homological hypotheses and derive the central conclusion. Given a surjection $F_0\to M$, the first syzygy module is its kernel. The decisive step is to use the universal property, exactness criterion, localization test, or finite-generation mechanism appropriate to the algebraic structure.
\end{example}

\begin{example}[Higher syzygies diagnostic]\label{ex:v24-worked-02}
Give a rigorous worked analysis of Higher syzygies. State the ring, module, finiteness, localization, or homological hypotheses and derive the central conclusion. The $n$th syzygy is the kernel at stage $n-1$ of a chosen free or projective resolution. The decisive step is to use the universal property, exactness criterion, localization test, or finite-generation mechanism appropriate to the algebraic structure.
\end{example}

\begin{example}[Dependence on presentation diagnostic]\label{ex:v24-worked-03}
Give a rigorous worked analysis of Dependence on presentation. State the ring, module, finiteness, localization, or homological hypotheses and derive the central conclusion. Syzygies depend on choices, but stable projective information is controlled up to adding free summands. The decisive step is to use the universal property, exactness criterion, localization test, or finite-generation mechanism appropriate to the algebraic structure.
\end{example}

\begin{example}[Presentations diagnostic]\label{ex:v24-worked-04}
Give a rigorous worked analysis of Presentations. State the ring, module, finiteness, localization, or homological hypotheses and derive the central conclusion. A finite presentation gives a finitely generated first syzygy. The decisive step is to use the universal property, exactness criterion, localization test, or finite-generation mechanism appropriate to the algebraic structure.
\end{example}

\section{Solved dossiers}
Each dossier is a canonical solved problem. The provenance ledger records which retained corpus rules guided its topic and which dossiers were added freshly to complete the chapter.

\begin{problem}[First syzygy diagnostic]\label{prob:v24-dossier-01}
Give a rigorous worked analysis of First syzygy. State the ring, module, finiteness, localization, or homological hypotheses and derive the central conclusion.
\end{problem}
\begin{solution}
Given a surjection $F_0\to M$, the first syzygy module is its kernel. The decisive step is to use the universal property, exactness criterion, localization test, or finite-generation mechanism appropriate to the algebraic structure.
\end{solution}

\begin{problem}[Higher syzygies diagnostic]\label{prob:v24-dossier-02}
Give a rigorous worked analysis of Higher syzygies. State the ring, module, finiteness, localization, or homological hypotheses and derive the central conclusion.
\end{problem}
\begin{solution}
The $n$th syzygy is the kernel at stage $n-1$ of a chosen free or projective resolution. The decisive step is to use the universal property, exactness criterion, localization test, or finite-generation mechanism appropriate to the algebraic structure.
\end{solution}

\begin{problem}[Dependence on presentation diagnostic]\label{prob:v24-dossier-03}
Give a rigorous worked analysis of Dependence on presentation. State the ring, module, finiteness, localization, or homological hypotheses and derive the central conclusion.
\end{problem}
\begin{solution}
Syzygies depend on choices, but stable projective information is controlled up to adding free summands. The decisive step is to use the universal property, exactness criterion, localization test, or finite-generation mechanism appropriate to the algebraic structure.
\end{solution}

\begin{problem}[Presentations diagnostic]\label{prob:v24-dossier-04}
Give a rigorous worked analysis of Presentations. State the ring, module, finiteness, localization, or homological hypotheses and derive the central conclusion.
\end{problem}
\begin{solution}
A finite presentation gives a finitely generated first syzygy. The decisive step is to use the universal property, exactness criterion, localization test, or finite-generation mechanism appropriate to the algebraic structure.
\end{solution}

\begin{problem}[Noetherian setting diagnostic]\label{prob:v24-dossier-05}
Give a rigorous worked analysis of Noetherian setting. State the ring, module, finiteness, localization, or homological hypotheses and derive the central conclusion.
\end{problem}
\begin{solution}
Over a Noetherian ring, syzygies of finite modules in finite free presentations remain finitely generated. The decisive step is to use the universal property, exactness criterion, localization test, or finite-generation mechanism appropriate to the algebraic structure.
\end{solution}

\begin{problem}[Dimension shifting diagnostic]\label{prob:v24-dossier-06}
Give a rigorous worked analysis of Dimension shifting. State the ring, module, finiteness, localization, or homological hypotheses and derive the central conclusion.
\end{problem}
\begin{solution}
Higher Tor and Ext groups can be expressed through lower groups of syzygies. The decisive step is to use the universal property, exactness criterion, localization test, or finite-generation mechanism appropriate to the algebraic structure.
\end{solution}

\begin{problem}[Regular sequences diagnostic]\label{prob:v24-dossier-07}
Give a rigorous worked analysis of Regular sequences. State the ring, module, finiteness, localization, or homological hypotheses and derive the central conclusion.
\end{problem}
\begin{solution}
Koszul-type relations provide canonical syzygies for complete-intersection situations. The decisive step is to use the universal property, exactness criterion, localization test, or finite-generation mechanism appropriate to the algebraic structure.
\end{solution}

\begin{problem}[Hilbert syzygy preview diagnostic]\label{prob:v24-dossier-08}
Give a rigorous worked analysis of Hilbert syzygy preview. State the ring, module, finiteness, localization, or homological hypotheses and derive the central conclusion.
\end{problem}
\begin{solution}
Over polynomial rings over fields, finite modules have bounded projective resolution length. The decisive step is to use the universal property, exactness criterion, localization test, or finite-generation mechanism appropriate to the algebraic structure.
\end{solution}

\begin{problem}[Syzygies over Noetherian rings are finite proof dossier]\label{prob:v24-dossier-09}
Prove the following structural statement and identify the hypothesis that makes the conclusion work: If $A$ is Noetherian and $M$ is finitely generated, then the kernel of a map $A^n\to M$ is finitely generated.
\end{problem}
\begin{solution}
The finite free module $A^n$ is Noetherian, so every submodule, including the kernel, is finitely generated. This proof isolates the reusable algebraic mechanism behind the result.
\end{solution}

\begin{problem}[Schanuel lemma proof dossier]\label{prob:v24-dossier-10}
Prove the following structural statement and identify the hypothesis that makes the conclusion work: If $0\to K\to P\to M\to0$ and $0\to K'\to P'\to M\to0$ have projective $P,P'$, then $K\oplus P'\cong K'\oplus P$.
\end{problem}
\begin{solution}
Form the pullback of $P\to M$ and $P'\to M$. Its two projections fit into short exact sequences that split because $P$ and $P'$ are projective, yielding the two decompositions. This proof isolates the reusable algebraic mechanism behind the result.
\end{solution}

\begin{problem}[Dimension shifting for Tor proof dossier]\label{prob:v24-dossier-11}
Prove the following structural statement and identify the hypothesis that makes the conclusion work: If $0\to\Omega M\to P\to M\to0$ with $P$ projective, then $\operatorname{Tor}_{n}^A(M,N)\cong\operatorname{Tor}_{n-1}^A(\Omega M,N)$ for $n\ge2$.
\end{problem}
\begin{solution}
Apply the long exact Tor sequence. Higher Tor of the projective module $P$ vanishes, so adjacent terms give the stated isomorphism. This proof isolates the reusable algebraic mechanism behind the result.
\end{solution}

\begin{problem}[Chapter synthesis]\label{prob:v24-dossier-12}
Connect the definitions and principal structural theorems of this chapter into one reusable commutative-algebra strategy.
\end{problem}
\begin{solution}
Syzygies are relations among generators, then relations among relations, and so on. Iterating syzygies turns presentations into resolutions and exposes hidden structure in modules. The recurring strategy is to encode the algebraic object by kernels, quotients, localization, finite presentation, or a resolution, apply the corresponding universal or exactness principle, and then recover the desired global or local conclusion.
\end{solution}

\section{Exercises with complete solutions}

\begin{exercise}\label{exr:v24-01}
State and apply the central principle behind First syzygy. Give one concrete algebraic consequence.
\end{exercise}
\begin{hint}
Given generators of M, place them in a surjection \(F_0\to M\); the kernel consists exactly of the first syzygies.
\end{hint}
\begin{solution}
Given a surjection $F_0\to M$, the first syzygy module is its kernel. The same principle can then be reused in later localization, support, base-change, resolution, Tor, or Ext arguments.
\end{solution}

\begin{exercise}\label{exr:v24-02}
State and apply the central principle behind Higher syzygies. Give one concrete algebraic consequence.
\end{exercise}
\begin{hint}
Write each relation as a coordinate vector in \(F_0\); those vectors generate the syzygy module.
\end{hint}
\begin{solution}
The $n$th syzygy is the kernel at stage $n-1$ of a chosen free or projective resolution. The same principle can then be reused in later localization, support, base-change, resolution, Tor, or Ext arguments.
\end{solution}

\begin{exercise}\label{exr:v24-03}
State and apply the central principle behind Dependence on presentation. Give one concrete algebraic consequence.
\end{exercise}
\begin{hint}
To compute higher syzygies, repeat the kernel computation for the map from a free module onto the previous syzygy module.
\end{hint}
\begin{solution}
Syzygies depend on choices, but stable projective information is controlled up to adding free summands. The same principle can then be reused in later localization, support, base-change, resolution, Tor, or Ext arguments.
\end{solution}

\begin{exercise}\label{exr:v24-04}
State and apply the central principle behind Presentations. Give one concrete algebraic consequence.
\end{exercise}
\begin{hint}
A presentation matrix encodes relations in its columns or rows according to convention; fix that convention before computing kernels.
\end{hint}
\begin{solution}
A finite presentation gives a finitely generated first syzygy. The same principle can then be reused in later localization, support, base-change, resolution, Tor, or Ext arguments.
\end{solution}

\begin{exercise}\label{exr:v24-05}
State and apply the central principle behind Noetherian setting. Give one concrete algebraic consequence.
\end{exercise}
\begin{hint}
Changing generators changes the presentation matrix by free-module operations but not the isomorphism class of M.
\end{hint}
\begin{solution}
Over a Noetherian ring, syzygies of finite modules in finite free presentations remain finitely generated. The same principle can then be reused in later localization, support, base-change, resolution, Tor, or Ext arguments.
\end{solution}

\begin{exercise}\label{exr:v24-06}
State and apply the central principle behind Dimension shifting. Give one concrete algebraic consequence.
\end{exercise}
\begin{hint}
Remove a redundant relation only after proving it lies in the submodule generated by the others.
\end{hint}
\begin{solution}
Higher Tor and Ext groups can be expressed through lower groups of syzygies. The same principle can then be reused in later localization, support, base-change, resolution, Tor, or Ext arguments.
\end{solution}

\begin{exercise}\label{exr:v24-07}
State and apply the central principle behind Regular sequences. Give one concrete algebraic consequence.
\end{exercise}
\begin{hint}
Syzygies are the data needed for the next differential in a free resolution; verify consecutive maps compose to zero.
\end{hint}
\begin{solution}
Koszul-type relations provide canonical syzygies for complete-intersection situations. The same principle can then be reused in later localization, support, base-change, resolution, Tor, or Ext arguments.
\end{solution}

\begin{exercise}\label{exr:v24-08}
State and apply the central principle behind Hilbert syzygy preview. Give one concrete algebraic consequence.
\end{exercise}
\begin{hint}
In graded settings, track degrees of generators and relations so the syzygy maps are homogeneous.
\end{hint}
\begin{solution}
Over polynomial rings over fields, finite modules have bounded projective resolution length. The same principle can then be reused in later localization, support, base-change, resolution, Tor, or Ext arguments.
\end{solution}

\section*{Chapter summary}
The chapter now contributes a canonical solved layer to the progression from rings and modules through localization, Noetherian methods, integral dependence, and derived functors.
% BEGIN VOL05-EXPANSION V24-exercises-01
\section{Graded supplementary exercises}
These exercises extend the protected chapter with a balanced set of computations, proofs, hypothesis tests, applications, and challenges.

\subsection*{Standard computations}

\begin{exercise}[Basic relation]\label{exr:v24-expand-01}
Give a syzygy between $x$ and $y$.
\end{exercise}
\begin{hint}
Use kernels of presentation maps and relations among generators.
\end{hint}
\begin{solution}
$(-y,x)$.
\end{solution}

\begin{exercise}[Kernel definition]\label{exr:v24-expand-02}
What is the first syzygy module in a presentation $F_1\to F_0\to M$?
\end{exercise}
\begin{hint}
Use kernels of presentation maps and relations among generators.
\end{hint}
\begin{solution}
The kernel of $F_0\to M$, equivalently the image from $F_1$ in an exact resolution.
\end{solution}

\begin{exercise}[Free module]\label{exr:v24-expand-03}
What are the positive syzygies of a free module in a length-zero resolution?
\end{exercise}
\begin{hint}
Use kernels of presentation maps and relations among generators.
\end{hint}
\begin{solution}
They vanish.
\end{solution}

\begin{exercise}[Cyclic quotient]\label{exr:v24-expand-04}
What is the first syzygy of $R/(f)$ in the standard presentation?
\end{exercise}
\begin{hint}
Use kernels of presentation maps and relations among generators.
\end{hint}
\begin{solution}
The principal submodule $fR\cong R$ when $f$ is regular.
\end{solution}

\begin{exercise}[Iterated kernel]\label{exr:v24-expand-05}
How is the second syzygy obtained?
\end{exercise}
\begin{hint}
Use kernels of presentation maps and relations among generators.
\end{hint}
\begin{solution}
As the kernel of a free map presenting the first syzygy.
\end{solution}

\subsection*{Proofs}

\begin{exercise}[Syzygies over Noetherian rings are finite proof]\label{exr:v24-expand-06}
Prove: If $A$ is Noetherian and $M$ is finitely generated, then the kernel of a map $A^n\to M$ is finitely generated.
\end{exercise}
\begin{hint}
Start from the theorem hypotheses and identify the decisive algebraic construction.
\end{hint}
\begin{solution}
The finite free module $A^n$ is Noetherian, so every submodule, including the kernel, is finitely generated.
\end{solution}

\begin{exercise}[Schanuel lemma proof]\label{exr:v24-expand-07}
Prove: If $0\to K\to P\to M\to0$ and $0\to K'\to P'\to M\to0$ have projective $P,P'$, then $K\oplus P'\cong K'\oplus P$.
\end{exercise}
\begin{hint}
Start from the theorem hypotheses and identify the decisive algebraic construction.
\end{hint}
\begin{solution}
Form the pullback of $P\to M$ and $P'\to M$. Its two projections fit into short exact sequences that split because $P$ and $P'$ are projective, yielding the two decompositions.
\end{solution}

\begin{exercise}[Dimension shifting for Tor proof]\label{exr:v24-expand-08}
Prove: If $0\to\Omega M\to P\to M\to0$ with $P$ projective, then $\operatorname{Tor}_{n}^A(M,N)\cong\operatorname{Tor}_{n-1}^A(\Omega M,N)$ for $n\ge2$.
\end{exercise}
\begin{hint}
Start from the theorem hypotheses and identify the decisive algebraic construction.
\end{hint}
\begin{solution}
Apply the long exact Tor sequence. Higher Tor of the projective module $P$ vanishes, so adjacent terms give the stated isomorphism.
\end{solution}

\begin{exercise}[Structural synthesis]\label{exr:v24-expand-09}
Prove a corollary that genuinely uses both Syzygies over Noetherian rings are finite and Schanuel lemma.
\end{exercise}
\begin{hint}
Apply the first theorem to obtain its structural description, then verify the second theorem's hypotheses.
\end{hint}
\begin{solution}
A correct proof explicitly invokes both results, verifies the common hypotheses, and derives a new consequence rather than merely restating either theorem.
\end{solution}

\subsection*{Counterexamples and hypothesis tests}

\begin{exercise}[No relations]\label{exr:v24-expand-10}
Test the claim: Every generating set has zero syzygy module.
\end{exercise}
\begin{hint}
Use the smallest concrete ring, module, or resolution that can decide the claim.
\end{hint}
\begin{solution}
False; relations are usually nontrivial.
\end{solution}

\begin{exercise}[Generator dependence]\label{exr:v24-expand-11}
Test the claim: The concrete syzygy module is independent of the chosen presentation on the nose.
\end{exercise}
\begin{hint}
Use the smallest concrete ring, module, or resolution that can decide the claim.
\end{hint}
\begin{solution}
False; stable isomorphism phenomena replace literal equality.
\end{solution}

\begin{exercise}[Finite syzygy]\label{exr:v24-expand-12}
Test the claim: Every syzygy over every ring is finitely generated.
\end{exercise}
\begin{hint}
Use the smallest concrete ring, module, or resolution that can decide the claim.
\end{hint}
\begin{solution}
False without Noetherian or finiteness hypotheses.
\end{solution}

\subsection*{Applications and investigations}

\begin{exercise}[Equation relations]\label{exr:v24-expand-13}
What does a syzygy encode computationally?
\end{exercise}
\begin{hint}
Translate the question into kernels of presentation maps and relations among generators.
\end{hint}
\begin{solution}
It records relations among chosen generators.
\end{solution}

\begin{exercise}[Resolution building]\label{exr:v24-expand-14}
How do syzygies build a free resolution?
\end{exercise}
\begin{hint}
Translate the question into kernels of presentation maps and relations among generators.
\end{hint}
\begin{solution}
Resolve the relation module repeatedly by new free modules.
\end{solution}

\subsection*{Challenge problems}

\begin{exercise}[Local-global synthesis]\label{exr:v24-expand-15}
Connect 'First syzygy' with 'Computational matrices' in one rigorous argument.
\end{exercise}
\begin{hint}
State the relevant ring map, module map, or complex before applying the structural theorem.
\end{hint}
\begin{solution}
The opening idea is: Given a surjection $F_0\to M$, the first syzygy module is its kernel. The later viewpoint is: Syzygies are kernels of matrices and can be studied algorithmically. A complete solution identifies the structural theorem that links these descriptions and checks its hypotheses.
\end{solution}

\begin{exercise}[Hypothesis audit]\label{exr:v24-expand-16}
Choose one theorem from the chapter, remove one essential hypothesis, and exhibit or explain a concrete failure.
\end{exercise}
\begin{hint}
Use one of the chapter's explicit computations or counterexamples as the test case.
\end{hint}
\begin{solution}
The solution must name the removed hypothesis, show exactly which proof step breaks, and give a concrete algebraic object where the claimed conclusion fails.
\end{solution}

% END VOL05-EXPANSION V24-exercises-01
